\chapter{Sprint 44 --- Chuẩn hoá nhãn backend sang tiếng Anh, portal giữ tiếng Việt (2026-08-02)}

\section{Bối cảnh}

Backend Wujia đang là \textbf{hỗn hợp hai ngôn ngữ}: nhãn do Odoo core sinh ra thì tiếng Anh, nhãn do
19 module Wujia tự khai thì tiếng Việt. Người dùng backend là đội Ngô Gia (một nhóm nhỏ, có nghiệp vụ
ERP); người dùng portal là \textbf{$\sim$1500 chủ tiệm nhượng quyền} --- hai tập người dùng khác hẳn
nhau nhưng đang ăn chung một kho chuỗi.

Mục tiêu sprint: \textbf{màn quản trị đọc bằng tiếng Anh, portal Y NGUYÊN tiếng Việt}, và có file
\texttt{vi\_VN.po} đầy đủ để đổi ngôn ngữ được \emph{hai chiều} --- ai muốn xem backend tiếng Việt
vẫn bật lại được, thấy đúng chữ như trước sprint.

\section{Cách làm --- dịch tại chỗ, không sed toàn cục}

\texttt{scripts/ba\_spec/i18n\_sprint44/i18n\_tool.py} (dev-only, gitignored) quét \emph{đúng những
vị trí được coi là nhãn backend} rồi thay tại chỗ:

\begin{center}
\begin{tabular}{p{3.2cm}p{10.6cm}}
\hline
\textbf{Loại file} & \textbf{Vị trí được nhận là nhãn backend} \\
\hline
XML attr & \texttt{string=} \texttt{help=} \texttt{confirm=} \texttt{placeholder=},
  \texttt{<menuitem name=>} \\
XML node & \texttt{<field name="name">} của act\_window / menuitem / res.groups / ir.rule; text
  trong \texttt{<t t-name="card">}; text trong help html \\
PY & \texttt{string=} \texttt{help=} \texttt{\_description=} và nhãn trong \texttt{Selection([...])} \\
\hline
\end{tabular}
\end{center}

\textbf{Loại trừ tuyệt đối}: mọi template portal (\texttt{portal\_*.xml}, \texttt{mobile\_*.xml},
\texttt{pc\_*.xml}, \texttt{*sidenav*}, \texttt{layouts}, \texttt{templates}\ldots), thư mục
\texttt{data/} (dữ liệu nghiệp vụ hiện trên portal), mọi text \texttt{ValidationError} /
\texttt{UserError}, comment và docstring, thư mục \texttt{tests/}.

Kết quả đợt đầu: \textbf{636 chuỗi trên 81 file}, kèm \textbf{14 file \texttt{i18n/vi\_VN.po}} sinh
từ \texttt{trans\_export} thật (không gõ tay) để chiều ngược lại luôn khớp.

\section{Ba việc phát sinh --- và vì sao chúng là việc thật}

\subsection{Va chạm động từ / trạng thái (13 sửa)}

Script dịch hàng loạt ánh xạ \emph{một} chuỗi Việt sang \emph{một} chuỗi Anh. Nhưng tiếng Việt dùng
chung một từ cho hai vai: \emph{``Từ chối''} vừa là \textbf{nút bấm} vừa là \textbf{trạng thái}. Sau
khi dịch, form phiếu đổi trả hiện \texttt{Reject} (nút) cạnh \texttt{Reject} (bậc statusbar) ---
người dùng không phân biệt được cái nào bấm được.

Nguyên tắc chốt: \textbf{nút $=$ động từ, state/filter $=$ tính từ}.

\begin{center}
\begin{tabular}{p{5.6cm}p{3.4cm}p{4.4cm}}
\hline
\textbf{Chỗ} & \textbf{Trước} & \textbf{Sau} \\
\hline
STATE của return / info-request / exam & \texttt{Reject} & \texttt{Rejected} \\
state \texttt{archived} của notification / fleet pricelist & \texttt{Archive} & \texttt{Archived} \\
filter tương ứng (4 chỗ) & \texttt{Reject} / \texttt{Archive} & \texttt{Rejected} / \texttt{Archived} \\
notebook page lý do & \texttt{Reject} & \texttt{Rejection reason} \\
nút \texttt{action\_lock\_portal} & \texttt{Portal locked} & \texttt{Lock portal} \\
nút \texttt{action\_delivery\_assign} & \texttt{Vehicle assigned} & \texttt{Assign vehicle} \\
\hline
\end{tabular}
\end{center}

Sáu nút \texttt{Reject}/\texttt{Archive} thật (hành động) \textbf{giữ nguyên} --- chúng vốn đã đúng.

Điểm quan trọng: \texttt{.po} chỉ cấm \emph{hai msgid trùng nhau}, không cấm hai msgid trả cùng một
msgstr. Nên \texttt{Reject} và \texttt{Rejected} cùng ra \emph{``Từ chối''}, \texttt{Archive} và
\texttt{Archived} cùng ra \emph{``Lưu trữ''}: người xem \texttt{vi\_VN} thấy chữ \textbf{y hệt trước
sprint}, còn người xem \texttt{en\_US} có hai từ phân biệt.

\subsection{Ba chỗ portal đọc thẳng nhãn của model}

Đây là loại lỗi nguy hiểm nhất của sprint: template portal \emph{không chứa} chữ tiếng Anh nào,
nhưng nó đọc \texttt{record.\_fields['x'].selection} --- mà selection vừa bị dịch. Portal đổi ngôn
ngữ theo backend một cách gián tiếp.

\begin{center}
\begin{tabular}{p{3.4cm}p{4.6cm}p{5.4cm}}
\hline
\textbf{Seam} & \textbf{Chỗ hở} & \textbf{Bịt bằng} \\
\hline
W2a trạng thái cửa hàng & 3 site QWeb đọc
  \texttt{franchise.\_fields['status']} & \texttt{FRANCHISE\_STATUS\_LABELS} trong
  \texttt{wujia\_portal\_base/controllers} \\
W2b mức độ ở popup chuông & \texttt{/portal/notification/recent} trả
  \texttt{priority\_label} compute từ model & \texttt{PORTAL\_PRIORITY\_LABELS}; \texttt{PC\_PRIORITY\_TAGS}
  đọc lại từ hằng này để còn 1 nguồn \\
W2c loại thông tin & controller đưa thẳng \texttt{REQUEST\_TYPE} ra dropdown &
  \texttt{REQUEST\_TYPE\_LABELS} + \texttt{REQUEST\_TYPE\_OPTIONS} \\
\hline
\end{tabular}
\end{center}

Pattern đã có sẵn trong repo (\texttt{BATCH\_STATUS\_LABELS} của
\texttt{wujia\_portal\_purchase\_history}, dựng ở Sprint 43): hằng nhãn Việt pin cứng trong
controller portal, kèm comment \emph{``key phải khớp \texttt{<CONST>} ở models/''}. Giữ đúng tinh
thần Sprint 42: \textbf{FE vẫn đọc nhãn do backend gửi}, không tự chế bảng nhãn --- chỉ đổi
\emph{nguồn} nhãn từ model sang hằng portal. \texttt{header\_bell\_badge.js} không sửa một dòng.

\subsection{Ba nhãn Việt \emph{không dấu} lọt lưới}

\texttt{i18n\_tool.py} lọc chuỗi Việt bằng \textbf{bảng ký tự có dấu}. Nên
\texttt{extract} trả về \emph{0 chuỗi còn lại} trong khi backend vẫn còn \texttt{Ca thi} (4 site),
\texttt{SL giao} và \texttt{\_description = 'Wujia Fleet Vehicle --- Xe'} --- tất cả đều ASCII thuần.
Phát hiện ra nhờ đi hết 31 menu backend ở \texttt{en\_US} và \emph{nhìn}: menu \emph{Ca thi} nằm giữa
một dàn menu tiếng Anh.

\begin{center}
\begin{tabular}{p{5.0cm}p{4.0cm}p{4.4cm}}
\hline
\textbf{Chỗ} & \textbf{en\_US} & \textbf{vi\_VN (giữ nguyên)} \\
\hline
menu + list + form + act\_window ca thi & \texttt{Time slot} & \emph{Ca thi} \\
\texttt{delivery\_qty} của wizard bù hàng & \texttt{Delivery qty} & \emph{SL giao} \\
\texttt{\_description} của \texttt{wujia.fleet.management} & \texttt{Wujia Fleet Vehicle} &
  \emph{Wujia Fleet Vehicle --- Xe} \\
\hline
\end{tabular}
\end{center}

\section{Bẫy lớn nhất của sprint --- \texttt{.pot} cũ nuốt cả file \texttt{.po}}

Sau khi nạp \texttt{.po} bằng \texttt{-{}-i18n-overwrite}, SQL đối chiếu vẫn báo \textbf{38 dòng có
\texttt{en\_US} nhưng \texttt{vi\_VN} rỗng} --- đúng 38 dòng đó lại \emph{có} msgstr đúng trong file
\texttt{.po}. File đúng, nạp xong, mà DB không đổi.

Nguyên nhân nằm ở \texttt{odoo/tools/translate.py} (\texttt{PoFileReader.\_\_init\_\_}): khi nạp
\texttt{<module>/i18n/vi\_VN.po}, Odoo \textbf{tự suy ra} \texttt{<module>.pot} cùng thư mục rồi
\texttt{merge} vào:

\begin{lstlisting}[style=python]
filename = path.parent.parent.name + '.pot'
pot_path = path.with_name(filename)
...
if pot_path:
    self.pofile.merge(polib.pofile(pot_path))
\end{lstlisting}

Ba module (\texttt{wujia\_sale}, \texttt{wujia\_portal\_knowledge},
\texttt{wujia\_portal\_support}) --- đúng ba module có sẵn \texttt{.pot} từ sprint trước --- vẫn giữ
\texttt{.pot} với msgid \textbf{tiếng Việt cũ}. \texttt{polib.merge} lấy \texttt{.pot} làm chuẩn, nên
mọi msgid tiếng Anh mới bị đánh dấu \texttt{obsolete}, và \texttt{\_\_iter\_\_} thì
\texttt{if entry.obsolete: continue}. File \texttt{.po} nạp vào nhưng \textbf{rỗng về mặt hiệu lực}.

Bằng chứng: \texttt{translation\_file\_reader} trên \texttt{wujia\_sale/i18n/vi\_VN.po} trả 85 dòng
với đúng \textbf{7 dòng không rỗng}, và \texttt{src} của chúng là các chuỗi Việt \emph{không còn tồn
tại trong file} (\emph{``Đơn từ Portal''}, \emph{``Cửa hàng nhượng quyền''}).

Cách sửa: \texttt{gen\_po.py} sinh lại \texttt{.pot} \emph{cùng lúc} với \texttt{.po} cho mọi module
có \texttt{.pot}. Sau đó số dòng thiếu tụt \textbf{38 $\rightarrow$ 0}.

\section{Kiểm chứng}

\begin{itemize}
  \item \textbf{Build}: 15 module, \texttt{-u \$MODS -{}-i18n-overwrite -{}-test-enable
    -{}-stop-after-init} $\Rightarrow$ \textbf{RC$=$0}, \textbf{49 post-test, 0 failed / 0 error},
    \texttt{Traceback} $=$ \textbf{0}. (49 test này thuộc \texttt{wujia\_portal\_debt} 14 $+$
    \texttt{wujia\_portal\_notification} 35; 13 module còn lại không có thư mục \texttt{tests/} ---
    nói \emph{``đã chạy 49 test''}, không nói \emph{``đã test 15 module''}.)
  \item \textbf{SQL hai ngôn ngữ} (\texttt{verify\_i18n.py}): soát \texttt{ir\_model\_fields},
    \texttt{ir\_model\_fields\_selection}, \texttt{ir\_ui\_menu}, \textbf{chỉ những dòng do 15 module
    Wujia sở hữu} (join \texttt{ir\_model\_data}) $\Rightarrow$ \textbf{0 dòng thiếu}, và 5 nhãn tách
    mới đều ra đúng chữ Việt đã chốt.
  \item \textbf{Đổi ngôn ngữ hai chiều} (\texttt{toggle\_lang\_check.py}): so \emph{token chính xác}
    trong header form. \texttt{en\_US} cho \texttt{Reject}(nút, phiếu \texttt{submitted}) và
    \texttt{Rejected}(state, phiếu \texttt{rejected}) tách bạch; \texttt{vi\_VN} cho \emph{``Từ
    chối''} cả hai. Notification và bảng giá xe hiện \emph{đồng thời} \texttt{Archive} $+$
    \texttt{Archived} --- đúng chỗ trước W1 hiện hai chữ y hệt.
  \item \textbf{Portal 391$\times$844} (\texttt{portal\_vi\_smoke.py}): 9 trang, chạy \emph{hai lượt}
    --- admin ở \texttt{vi\_VN} và ở \texttt{en\_US}. Lượt \texttt{en\_US} mới là lượt quan trọng vì
    đó là trạng thái sau sprint. Kết quả: 200 hết, \textbf{overflow ngang $=$ 0}, không rò token
    tiếng Anh. Ba seam soi bằng dữ liệu thật: popup chuông trả \emph{Cần làm / Quan trọng / Thông
    thường}; dropdown loại thông tin ra 7 nhãn Việt; badge trạng thái cửa hàng ra tiếng Việt.
  \item \textbf{Menu backend} (\texttt{menu\_walk.py}): \textbf{31/31} menu có action mở được ở
    \texttt{en\_US}, không 500 / OwlError / error dialog. \texttt{check\_kanban.py}: cả hai kanban
    (Hỗ Trợ, Kiến thức) vẫn render sau fix \texttt{kanban-box} $\rightarrow$ \texttt{card} của Odoo 19.
\end{itemize}

\section{Chưa làm --- và vì sao}

\begin{itemize}
  \item \textbf{11 finding ``dịch chưa sát nghĩa''} (Province/City, Severity, nhãn nút wizard\ldots)
    --- defer chờ BA chốt từ vựng. Nguồn: \texttt{review\_findings\_raw.txt}.
  \item \textbf{Nhãn \texttt{wujia.franchise.member.role}} (\emph{Chủ tiệm / Quản lý / Nhân viên})
    vẫn tiếng Việt: nó nằm trong \texttt{display\_name} đã store, đổi cần migration.
  \item \textbf{Menu vốn đã tiếng Anh từ trước} (\emph{Franchise Management}, \emph{Fleet
    Management}\ldots) \textbf{chưa có bản dịch \texttt{vi\_VN}} --- ai xem backend bằng
    \texttt{vi\_VN} sẽ thấy các menu đó bằng tiếng Anh. Đây là tồn có sẵn, không phải hồi quy.
  \item \textbf{Hai chỗ portal vốn đã tiếng Anh từ trước sprint}, đã đối chiếu \texttt{git show
    HEAD}: \texttt{portal\_support.xml:430} (\emph{Normal/Urgent}) và
    \texttt{portal\_franchise\_information.xml:248--249} (badge thành viên \emph{Active/Inactive},
    trong khi bản PC dòng 124 lại ghi \emph{``Đang hoạt động''} --- lệch sẵn). Không phải hồi quy
    sprint 44; ghi vào tồn để BA chốt.
  \item \textbf{UAT}: đã đổi \texttt{lang} của admin sang \texttt{en\_US} trên LOCAL; UAT chỉ đổi
    \textbf{sau khi deploy} code sprint 44 --- đổi trước thì màn quản trị lẫn lộn tiếng Anh của core
    với tiếng Việt của source cũ, còn rối hơn.
\end{itemize}

\section{Bài học}

\begin{itemize}
  \item \textbf{Một từ tiếng Việt có thể là hai từ tiếng Anh.} Dịch 1--1 làm hỏng chỗ mà tiếng Việt
    dùng chung từ cho nút và cho trạng thái. Quy tắc để lại: \emph{nút $=$ động từ, state/filter $=$
    tính từ}; và hai msgid trả cùng một msgstr là hợp lệ, nên chiều \texttt{vi\_VN} không mất gì.
  \item \textbf{Template không chứa chữ tiếng Anh vẫn có thể rò tiếng Anh.} Chỗ hở là
    \texttt{\_fields[...].selection}, không phải chữ trong file. Muốn portal độc lập ngôn ngữ thì
    phải \emph{pin nhãn ở controller portal}, không đọc nhờ nhãn của model.
  \item \textbf{Odoo tự merge \texttt{<module>.pot} khi nạp \texttt{.po}.} Một \texttt{.pot} cũ đủ
    làm cả file \texttt{.po} mất tác dụng \emph{trong im lặng} --- không lỗi, không cảnh báo, chỉ là
    DB không đổi. Sinh \texttt{.po} thì phải sinh \texttt{.pot} cùng lúc.
  \item \textbf{Bộ lọc theo dấu tiếng Việt bỏ sót tiếng Việt không dấu.} \emph{Ca thi}, \emph{SL
    giao}, \emph{Xe} đều là ASCII. \texttt{extract} báo 0 không có nghĩa là xong --- phải đi hết menu
    và \emph{nhìn}.
  \item \textbf{Assert bằng \texttt{in} là assert rỗng khi một chuỗi là con của chuỗi kia.}
    \texttt{"Reject" in "Rejected"} luôn True, nên bài test đầu tiên pass mà không chứng minh gì.
    Phải so token chính xác.
  \item \textbf{Overlay che trang làm mọi assert nội dung pass rỗng.} Admin là owner 3 cửa hàng và
    chưa chọn cửa hàng nào $\Rightarrow$ \texttt{\#wujiaStoreOverlay} phủ toàn trang;
    \texttt{innerText} chỉ trả text của overlay. Phải set sẵn cookie
    \texttt{wujia\_active\_franchise\_id} trước khi smoke portal.
  \item \textbf{\texttt{pkill -f 'odoo[-]bin'} chỉ an toàn khi dòng lệnh không chứa
    \texttt{odoo-bin} ở chỗ khác.} Gộp \texttt{pkill} chung một dòng với lệnh chạy
    \texttt{odoo-bin} thì shell tự giết chính nó (exit 144). Tách thành hai lệnh riêng.
  \item \textbf{\texttt{odoo.conf} có \texttt{logfile=} nên redirect stdout bắt được gần như không
    gì.} Luôn thêm \texttt{-{}-logfile=} khi cần đọc log build.
\end{itemize}
